\section{The Dirichlet problem on a periodic fibre}\label{app:dirichlet}

The energy argument of Section \ref{sec:energy} uses a potential that solves a quaternionic
Monge--Amp\`ere equation on each spatial fibre, is smooth up to a boundary that moves with the
parameters, and is compared with explicit functions on the fibre and on small balls inside it. This
appendix proves the existence, comparison and smooth dependence that are needed; the quantitative
bounds in the wavelength come from the comparison functions of Section \ref{sec:fiber}. The argument
is that of \cite[Appendix A]{OAI26}, written out in full. Alesker solved the Dirichlet problem on
quaternionic balls \cite[Theorem 1.4]{Ale03b} and Zhu on bounded domains of $\HH^m$ carrying a smooth
subsolution \cite[Theorem 1.1]{Zhu16}. Our fibres have periodic real directions and only three real
coordinates per variable; a fourth, periodic coordinate is adjoined to provide the quaternionic
rotations used for the mixed derivatives at the boundary.

\subsection{Setting and statement}
Let $m\ge1$, let $\Lambda\subset\R^m$ be a lattice of rank $m$, let $V\subset\Ck^m\cong\R^{2m}$ be a bounded
domain with smooth boundary, and put $\Omega=(\R^m/\Lambda)\times V$. In the notation of Section
\ref{ssec:hessian} with three slots, $z_a=x_a+P_a\ii$, the periodic variables $x_a$ occupy the real
slot and $P_a=p_a+\kk q_a\in\Ck$ the slots $\ii,\jj$; $H$ is the three-slot Hessian \eqref{eq:GH}. A
function on $V$ is regarded as a function on $\Omega$ independent of $x$.

\begin{theorem}[{\cite[Theorem A.1]{OAI26}}]\label{thm:dirichlet}
Let $r$ be smooth on a neighbourhood of $\overline V$, with $r<0$ in $V$, $r=0$ and $dr\ne0$ on
$\partial V$, and $H(r)>0$ on $\overline\Omega$. Let $B_0,g\in C^\infty(\overline\Omega)$ be real with
\[
H(B_0)\ge0,\qquad g>\det H(B_0)\qquad\text{on }\overline\Omega .
\]
Then there is a unique $u\in C^\infty(\overline\Omega)$ with
\begin{equation}\label{eq:dirichletA}
\det H(B_0+u)=g,\qquad H(B_0+u)>0\quad\text{on }\overline\Omega,\qquad u|_{\partial\Omega}=0,
\end{equation}
and there are constants $b\ge b'>0$ with
\begin{equation}\label{eq:barrier}
b\,r\le u\le b'\,r .
\end{equation}
If $r,B_0,g$ and the domain depend smoothly on finitely many parameters, the constants $b,b'$ can be
chosen locally uniformly in the parameters, and the solution depends smoothly on the parameters up to
the boundary after any smooth local identification of the domains with a fixed domain.
\end{theorem}

Here $C^\infty(\overline\Omega)$ denotes functions all of whose derivatives extend continuously to
$\overline\Omega$, and similarly for the H\"older spaces $C^{k,\alpha}(\overline\Omega)$; $C^{k,\alpha}_0$ is
the subspace of functions vanishing on $\partial\Omega$. The \emph{smooth dependence} means: if
$\Theta_\tau:\overline V_{\tau_0}\to\overline V_\tau$ are diffeomorphisms depending smoothly on $(\tau,y)$, then
$(\tau,x,y)\mapsto u_\tau(x,\Theta_\tau(y))$ is smooth up to the boundary.

\subsection{Matrix facts}
We use the following consequences of Lemma \ref{lem:moore}. Let $Q,Q_1,Q_2$ be positive definite
Hermitian quaternionic $m\times m$ matrices and $B$ Hermitian.
\begin{enumerate}[label=(M\arabic*)]
\item If $A\ge0$ and $B\ge0$ then $A:B\ge0$. Indeed $A=UDU^*$ with $U$ unitary and $D\ge0$ diagonal, and
$A:B=\Re\tr(D^{1/2}U^*BUD^{1/2})$ by cyclicity of $\Re\tr$, a sum of the nonnegative numbers
$e_i^*D^{1/2}U^*BUD^{1/2}e_i$.
\item $\frac{d}{ds}\log\det(Q+sB)|_{s=0}=Q^{-1}:B$, by \eqref{eq:cofactor} and \eqref{eq:cofinv}. Since
$Q_\theta=(1-\theta)Q_2+\theta Q_1>0$ for $0\le\theta\le1$,
\begin{equation}\label{eq:logdiff}
\log\det Q_1-\log\det Q_2=\Bigl(\int_0^1Q_\theta^{-1}\,d\theta\Bigr):(Q_1-Q_2).
\end{equation}
\item Put $S=\Re\tr Q^{-1}$, the sum of the reciprocals of the eigenvalues of $Q$. By the
arithmetic--geometric mean inequality $S\ge m(\det Q)^{-1/m}$. If $-\beta I\le B\le\beta I$ then
$|Q^{-1}:B|\le\beta S$, by (M1) applied to $\beta I\pm B$.
\item If $A\ge0$ and $P\ge cI$ with $c>0$, then $\det(A+P)>\det A$. Indeed the real representation
\eqref{eq:realrep} turns $A$ and $A+P$ into symmetric real matrices with $(A+P)_{\R}\ge A_{\R}+cI$, so by the
minimax characterisation each eigenvalue of $(A+P)_{\R}$ exceeds the corresponding eigenvalue of $A_{\R}$
by at least $c$, and \eqref{eq:realrep} gives $\det(A+P)^4>\det(A)^4\ge0$.
\end{enumerate}

\subsection{The linearised operator and comparison}
Work with $k\in\{3,4\}$ slots and $m$ coordinate indices; the real dimension is $N=km$, and $H$ is the
$k$-slot Hessian. For a real covector $\xi\in\R^N$ write $\hat\xi=(\sum_\alpha e_\alpha\xi_{a\alpha})_a\in\HH^m$;
then $|\hat\xi|=|\xi|$. For a real symmetric $N\times N$ matrix $M$ let $H(M)$ be the Hermitian matrix with
$H(D^2v)=H(v)$; it is linear in $M$, and $H(\xi\xi^{\rm t})=\hat\xi\hat\xi^*$ by Lemma \ref{lem:hessian}. If
$Q>0$ is a Hermitian matrix field, the operator
\[
L_Qw=Q^{-1}:H(w)
\]
is a real linear operator $\sum_{i,j}a^{ij}\partial_i\partial_jw$ with no lower-order terms, whose coefficient
matrix is defined by $a:M=Q^{-1}:H(M)$. For $M=\xi\xi^{\rm t}$ this gives
$a(\xi,\xi)=\hat\xi^*Q^{-1}\hat\xi$, so
\begin{equation}\label{eq:ellipticity}
\lambda_{\max}(Q)^{-1}|\xi|^2\le a(\xi,\xi)\le\lambda_{\min}(Q)^{-1}|\xi|^2 .
\end{equation}
If $\psi$ is a $C^2$ function with $Q=H(\psi)>0$ we write $L_\psi=L_{H(\psi)}$, the derivative of
$\psi\mapsto\log\det H(\psi)$ by (M2).

In what follows a \emph{periodic domain} is an open set $U\subset(\R^{m}/\Lambda)\times W$, where $W$ is
bounded in the remaining real coordinates and $\Lambda$ is a lattice (possibly of lower rank, possibly
acting on more coordinates), so that $\overline U$ is compact and at least one real coordinate $y$ is
a bounded function on $U$. The domain $\Omega$, its four-slot extension below, and balls in a
fibre are periodic domains.

\begin{lemma}[Maximum principle]\label{lem:maxprinciple}
Let $U$ be a periodic domain and $L=\sum a^{ij}\partial_i\partial_j$ with $a$ continuous and positive
definite in $U$. If $f\in C^2(U)\cap C(\overline U)$ satisfies $Lf\ge0$ in $U$, then
$\max_{\overline U}f=\max_{\partial U}f$.
\end{lemma}

\begin{proof}
At an interior maximum $D^2f\le0$, so $Lf=a:D^2f\le0$; hence a function with $Lf>0$ has no interior
maximum. Let $y$ be a bounded real coordinate and $h=e^{y}$. Then $Lh=a^{yy}e^y>0$, so
for $\eps>0$ the function $f+\eps h$ satisfies $L(f+\eps h)>0$ and attains its maximum over the
compact set $\overline U$ on $\partial U$. Letting $\eps\to0$ gives the claim.
\end{proof}

\begin{lemma}[Comparison]\label{lem:comparison2}
Let $U$ be a periodic domain in $k$ slots, and $\psi_1,\psi_2\in C^2(U)\cap C(\overline U)$ with
$H(\psi_i)>0$ in $U$. If $\det H(\psi_1)\ge\det H(\psi_2)$ in $U$ and $\psi_1\le\psi_2$ on $\partial U$,
then $\psi_1\le\psi_2$ in $U$. In particular \eqref{eq:dirichletA} has at most one solution in
$C^2(\Omega)\cap C(\overline\Omega)$.
\end{lemma}

\begin{proof}
By \eqref{eq:logdiff}, $0\le\log\det H(\psi_1)-\log\det H(\psi_2)=L_{\bar Q}(\psi_1-\psi_2)$ with
$\bar Q^{-1}=\int_0^1Q_\theta^{-1}d\theta$, $Q_\theta=(1-\theta)H(\psi_2)+\theta H(\psi_1)$. The coefficients
are continuous and positive definite in $U$ by \eqref{eq:ellipticity}, and Lemma \ref{lem:maxprinciple}
gives $\psi_1-\psi_2\le\max_{\partial U}(\psi_1-\psi_2)\le0$. For uniqueness apply this twice with
$\psi_i=B_0+u_i$.
\end{proof}

\subsection{The four-slot extension}
Let $\tilde\Omega=\Omega\times(\R/\Z)^m$, with the new coordinates $z_{a3}$ in the $\kk$ slots, and extend
$r,B_0,g$ to $\tilde\Omega$ as functions independent of them. For a function independent of the $\kk$
slots the four-slot $G$ and $H$ of \eqref{eq:GH} agree with the three-slot ones, since every term
containing $\partial_{a3}$ vanishes. Hence the hypotheses of Theorem \ref{thm:dirichlet} hold on
$\tilde\Omega$, and a solution on $\Omega$ extends to a solution on $\tilde\Omega$. Conversely, suppose
$\tilde u$ solves the four-slot problem on $\tilde\Omega$. For $c\in(\R/\Z)^m$ the translate
$\tilde u(\cdot+c)$ in the $\kk$ slots solves the same problem, since the data are invariant and
$H$ commutes with translations. By Lemma \ref{lem:comparison2} (with $k=4$), $\tilde u(\cdot+c)=\tilde u$,
so $\tilde u$ is independent of the $\kk$ slots and restricts to a solution on $\Omega$. Therefore it
suffices to prove Theorem \ref{thm:dirichlet} on $\tilde\Omega$, and the rest of the proof works there,
with four slots and $N=4m$.

The manifold $\overline{\tilde\Omega}$ is compact with smooth boundary
$\partial\tilde\Omega=(\R^m/\Lambda)\times\partial V\times(\R/\Z)^m$ and a flat metric, and each point has a
neighbourhood isometric to an open subset of $\R^{4m}$ in which the boundary is a smooth hypersurface.
We call such a neighbourhood a \emph{flat chart}. By compactness $\overline{\tilde\Omega}$ is covered by
finitely many flat charts of a fixed size. The local Schauder estimates \cite[Corollary 6.3 and Lemma
6.5]{GT01}, their higher-order forms, and the local $L^2$ regularity \cite[Theorems 8.8 and 8.12]{GT01}
are local statements and apply in such charts; summing over a finite cover gives the global estimates
of \cite[Theorem 6.6]{GT01} and their higher-order analogues on $\overline{\tilde\Omega}$, exactly as in
the proofs in \cite[Ch.~6]{GT01}.

\subsection{The linear Dirichlet problem}
\begin{lemma}\label{lem:lineardirichlet}
Let $Q>0$ be a smooth Hermitian field on $\overline{\tilde\Omega}$, $k\ge0$ and $0<\alpha<1$. Then
$L_Q:C^{k+2,\alpha}_0(\overline{\tilde\Omega})\to C^{k,\alpha}(\overline{\tilde\Omega})$ is an isomorphism.
\end{lemma}

\begin{proof}
For the Laplacian: by Poincar\'e's inequality along a bounded coordinate of $V$, the form
$\int\nabla v\cdot\nabla w$ is coercive on $H^1_0(\tilde\Omega)$, so for $f\in C^\infty(\overline{\tilde\Omega})$ the
Lax--Milgram theorem gives a weak solution $v\in H^1_0$ of $\Delta v=f$. The local regularity theorems,
interior and at the flat boundary pieces, give $v\in H^j$ for every $j$, hence
$v\in C^\infty(\overline{\tilde\Omega})$ by Sobolev embedding. For $f\in C^{k,\alpha}$ mollify in charts to get
smooth $f_\ell$ bounded in $C^{k,\alpha}$ and converging in $C^{k,\beta}$, $\beta<\alpha$. The estimate
below bounds the solutions $v_\ell$ in $C^{k+2,\alpha}$; they converge in $C^{k+2,\beta}$, and the limit
lies in $C^{k+2,\alpha}_0$ and solves $\Delta v=f$. Uniqueness follows from Lemma \ref{lem:maxprinciple}.

For $L_s=(1-s)\Delta+sL_Q$, $0\le s\le1$, the coefficients are smooth and uniformly elliptic,
uniformly in $s$, since $\Delta=L_I$ by Lemma \ref{lem:hessian}. Since $H(r)\ge c_rI$ on
$\overline{\tilde\Omega}$ for some $c_r>0$, (M1) gives $L_sr\ge c'>0$ uniformly in $s$. If
$v\in C^2_0$ and $|L_sv|\le K$, then $L_s(\pm v+\frac K{c'}r)\ge0$ and Lemma \ref{lem:maxprinciple} gives
$|v|\le\frac K{c'}|r|$. With the global Schauder estimate
$\|v\|_{C^{k+2,\alpha}}\le C(\|v\|_{C^0}+\|L_sv\|_{C^{k,\alpha}})$ this gives
$\|v\|_{C^{k+2,\alpha}}\le C\|L_sv\|_{C^{k,\alpha}}$ on $C^{k+2,\alpha}_0$, uniformly in $s$, and the method of
continuity \cite[Theorem 5.2]{GT01} transfers invertibility from $L_0=\Delta$ to $L_1=L_Q$.
\end{proof}

\subsection{The continuity path}
Write $\psi=B_0+u$, so that \eqref{eq:dirichletA} reads $\det H(\psi)=g$, $H(\psi)>0$,
$\psi=B_0$ on $\partial\tilde\Omega$. For $0\le t\le1$ put
\[
g_t=(1-t)\det H(B_0+r)+tg .
\]
At $t=0$ the function $u=r$ is a solution: $H(B_0+r)\ge H(r)>0$ and $r=0$ on the boundary. By (M4),
$\det H(B_0+r)>\det H(B_0)$, and $g>\det H(B_0)$ by hypothesis; by compactness there are
$\delta,c_0,C_0>0$ with
\begin{equation}\label{eq:gap}
g_t\ge\det H(B_0)+\delta,\qquad c_0\le g_t\le C_0\qquad\text{on }\overline{\tilde\Omega},\ 0\le t\le1 .
\end{equation}
Let $T$ be the set of $t\in[0,1]$ for which \eqref{eq:dirichletA} with $g_t$ in place of $g$ has a
solution $u_t\in C^\infty(\overline{\tilde\Omega})$. Then $0\in T$, and we show that $T$ is open and closed.

\emph{Openness.} Let $t_0\in T$, fix $\alpha\in(0,1)$ and $k\ge0$, and consider
\[
\mathcal F(t,u)=\log\det H(B_0+u)-\log g_t
\]
on the open set of $u\in C^{k+2,\alpha}_0$ with $H(B_0+u)>0$ on $\overline{\tilde\Omega}$, with values in
$C^{k,\alpha}$. It is a smooth map: $\log\det H(M)$ is a smooth function of the real Hessian $M$ on the
open set $H(M)>0$, and composition with a smooth function is a smooth map of H\"older spaces. Its
derivative in $u$ at $u_{t_0}$ is $L_{\psi_{t_0}}$, an isomorphism by Lemma
\ref{lem:lineardirichlet}. The implicit function theorem gives, for $t$ near $t_0$, a solution
$u_t\in C^{k+2,\alpha}_0$ depending smoothly on $t$. By Lemma \ref{lem:comparison2} the solutions obtained
for different $k$ coincide, so $u_t\in C^\infty(\overline{\tilde\Omega})$ and $T$ is open.

\emph{Closedness.} By the a priori estimates below, every $u_t$, $t\in T$, satisfies
$\|u_t\|_{C^j(\overline{\tilde\Omega})}\le C_j$ for every $j$, and $H(B_0+u_t)\ge c\,I$, with $C_j,c$
independent of $t$. If $t_\ell\in T$ and $t_\ell\to t$, the Arzel\`a--Ascoli theorem gives a subsequence of
$u_{t_\ell}$ converging in every $C^j$ to some $u\in C^\infty(\overline{\tilde\Omega})$, which vanishes on the
boundary, satisfies $H(B_0+u)\ge cI$ and $\det H(B_0+u)=g_t$. So $t\in T$.

Hence $T=[0,1]$ and $u=u_1$ solves \eqref{eq:dirichletA}. It remains to prove the a priori estimates.
For the rest of this appendix $u=u_t$ is a smooth solution for some $t\in T$, $\psi=B_0+u$,
$Q=H(\psi)$, $L=L_\psi$, $S=\Re\tr Q^{-1}$, and $C$ denotes constants independent of $t$ and $u$.
By (M3) and \eqref{eq:gap},
\begin{equation}\label{eq:Slow}
S\ge m\,g_t^{-1/m}\ge mC_0^{-1/m}=:s_0>0,\qquad Lr=Q^{-1}:H(r)\ge c_rS .
\end{equation}

\subsection{Barriers and the gradient}
\emph{The barriers \eqref{eq:barrier}.} As $b'\downarrow0$, $\det H(B_0+b'r)\to\det H(B_0)$ uniformly, so
by \eqref{eq:gap} we may fix $b'\in(0,1)$ with $\det H(B_0+b'r)<g_t$ for all $t$. Since
$\det H(B_0+br)\ge\det(bH(r))\ge(bc_r)^m$ by (M4) and monotonicity, we may fix $b>1$ with
$\det H(B_0+br)>C_0\ge g_t$. All three potentials $B_0+b'r$, $\psi$, $B_0+br$ equal $B_0$ on the boundary
and have positive definite Hessians, so Lemma \ref{lem:comparison2} gives $B_0+br\le\psi\le B_0+b'r$,
which is \eqref{eq:barrier}. Since $r$ is a defining function, near $\partial\tilde\Omega$ we can write
$u=\rho\,r$ with $\rho$ smooth up to the boundary (Section \ref{ssec:divergence}). On the collar
$r<0$, so \eqref{eq:barrier} gives $b'\le\rho\le b$, and by continuity also on the boundary. There
$\nabla u=\rho\nabla r$, so
\begin{equation}\label{eq:boundarygrad}
|\nabla u|\le b\,|\nabla r|\qquad\text{on }\partial\tilde\Omega .
\end{equation}

\emph{First and second derivatives along constant directions.} For real symmetric $M$ with $H(M)>0$
put $F(M)=\log\det H(M)$. It is concave, being the composition of the linear map $M\mapsto H(M)$ with
the concave function $\log\det$ (Lemma \ref{lem:moore}), and $DF(M)[\dot M]=H(M)^{-1}:H(\dot M)$ by
(M2). For a constant real unit vector $e$ write $\psi_e,\psi_{ee}$ for the directional derivatives.
Differentiating $F(D^2\psi)=\log g_t$ once and twice along $e$ gives
\begin{equation}\label{eq:diffeq}
L\psi_e=(\log g_t)_e,\qquad L\psi_{ee}=(\log g_t)_{ee}-D^2F(D^2\psi)[D^2\psi_e,D^2\psi_e]\ge(\log g_t)_{ee}.
\end{equation}
By \eqref{eq:Slow}, $L(\pm\psi_e+Kr)\ge-C+Kc_rs_0>0$ for a large constant $K$, and Lemma
\ref{lem:maxprinciple} gives $\pm\psi_e\le\max_{\partial\tilde\Omega}|\nabla\psi|+K|r|$. With
\eqref{eq:boundarygrad} this bounds $|\nabla\psi|$, hence $|\nabla u|$, by a constant $C$. In the same
way $L(\psi_{ee}+Kr)>0$ for large $K$, so
\begin{equation}\label{eq:hessreduce}
\psi_{ee}\le\max_{\partial\tilde\Omega}|D^2\psi|+C\qquad\text{in }\tilde\Omega,\ |e|=1 .
\end{equation}
It remains to bound $D^2\psi$ on the boundary; $B_0$ is smooth, so it suffices to bound $D^2u$ there.

\subsection{Second derivatives at the boundary}
Fix a boundary point and a flat chart centred at it, with coordinates $x\in\R^{4m}$, identified with
$\HH^m$ through the slots, so that the real inner product is $\Re(v^*w)$. Let $\nu$ be the outward unit
normal at $0$; then $\nabla r(0)=|\nabla r(0)|\nu$, since $r<0$ inside. All constants below are uniform
over boundary points, by compactness.

\emph{Tangential directions.} From $u=\rho r$,
$D^2u=\rho D^2r+\nabla\rho\otimes\nabla r+\nabla r\otimes\nabla\rho+rD^2\rho$, and at $0$ the last three terms
vanish on pairs of tangent vectors. Hence $|D^2u(0)(v,w)|=\rho|D^2r(0)(v,w)|\le C$ for unit tangent
$v,w$.

\emph{Mixed directions.} Fix a unit tangent vector $e$ at $0$. Let $P_{\rm tan}$ be the real orthogonal
projection onto the tangent space and put
\[
v=\frac{P_{\rm tan}D^2r(0)e}{|\nabla r(0)|},\qquad A=v\nu^*-\nu v^*-\nu(\nu^*v)\nu^* ,
\]
with vectors read as quaternionic columns. Since $v$ is tangent, $\Re(\nu^*v)=0$, so the quaternion
$\nu^*v$ satisfies $\overline{\nu^*v}=-\nu^*v$. Hence
$A^*=\nu v^*-v\nu^*+\nu(\nu^*v)\nu^*=-A$, and, using $\nu^*\nu=1$ and $v^*\nu+\nu^*v=2\Re(\nu^*v)=0$,
\[
A\nu=v-\nu(v^*\nu)-\nu(\nu^*v)=v .
\]
Consider the affine vector field $X(x)=e+Ax$. Its flow is $\phi_s(x)=e^{sA}x+c_s$ with
$c_s=\int_0^se^{\sigma A}e\,d\sigma$, and $U_s=e^{sA}$ is quaternionic unitary since $A^*=-A$.

First, $Xr$ vanishes to second order along the boundary at $0$. Indeed $Xr(0)=\nabla r(0)\cdot e=0$, and
for a tangent vector $w$,
\begin{align*}
\partial_w(Xr)(0)&=D^2r(0)(w,e)+\nabla r(0)\cdot(Aw)
=|\nabla r(0)|\bigl(v\cdot w+\Re(\nu^*Aw)\bigr)\\
&=|\nabla r(0)|\bigl(v\cdot w-\Re((A\nu)^*w)\bigr)=0,
\end{align*}
because $D^2r(0)(w,e)=w\cdot P_{\rm tan}D^2r(0)e$ and $\Re(\nu^*Aw)=\Re((A^*\nu)^*w)=-\Re((A\nu)^*w)$.
So the restriction of $Xr$ to the boundary hypersurface vanishes at $0$ with its differential, and
$|Xr(x)|\le C|x|^2$ for boundary points $x$ in the chart. On the boundary $Xu=\rho\,Xr+r\,X\rho=\rho\,Xr$,
so
\begin{equation}\label{eq:Xuboundary}
|Xu(x)|\le C|x|^2\qquad(x\in\partial\tilde\Omega,\ |x|\le\varrho_0),
\end{equation}
for a fixed chart radius $\varrho_0$.

Secondly, $L(X\psi)=X\log g_t$. By the covariance in Lemma \ref{lem:hessian} and its invariance under
translations, $H(\psi\circ\phi_s)(x)=U_s^*H(\psi)(\phi_s(x))U_s$, and the Moore determinant is invariant
under unitary conjugation, since the eigenvalues are. Hence
$\log\det H(\psi\circ\phi_s)(x)=\log g_t(\phi_s(x))$ for $x$ in the chart and small $s$, and
differentiation at $s=0$, using (M2) and $\frac d{ds}\psi\circ\phi_s|_{s=0}=X\psi$, gives
$L(X\psi)=X\log g_t$. Since $X B_0$ is smooth with $|H(XB_0)|\le C$, (M3) and \eqref{eq:Slow} give
\begin{equation}\label{eq:LXu}
|L(Xu)|\le|X\log g_t|+|L(XB_0)|\le C+CS\le CS .
\end{equation}

Let $U=\{|x|<\varrho_0\}\cap\tilde\Omega$ and $w=K_1|x|^2-K_2r$. On the part of $\partial U$ in the
boundary of $\tilde\Omega$, $w=K_1|x|^2\ge|Xu|$ by \eqref{eq:Xuboundary} once $K_1\ge C$. On the part
where $|x|=\varrho_0$, $w\ge K_1\varrho_0^2\ge|X||\nabla u|\ge|Xu|$ once $K_1$ is large, by the gradient
bound. In four slots $H(|x|^2)=8I$, since $|x|^2=\sum_{a,\alpha}z_{a\alpha}^2$ has real Hessian $2I$ and
each diagonal entry collects $\sum_\alpha e_\alpha\overline{e_\alpha}\cdot2=8$. So by \eqref{eq:Slow} and
\eqref{eq:LXu},
\[
L(w\pm Xu)\le8K_1S-K_2c_rS+CS\le0
\]
once $K_2$ is large. Lemma \ref{lem:maxprinciple} applied to $-(w\pm Xu)$ gives $|Xu|\le w$ in $U$.
At $0$ both $w$ and $Xu=\rho Xr$ vanish. Along the inward normal, $x=-s\nu$ with small $s>0$, we get
$|Xu(-s\nu)|\le K_1s^2-K_2r(-s\nu)\le Cs$, so $|\partial_\nu(Xu)(0)|\le C$. Now
\[
\partial_\nu(Xu)(0)=D^2u(0)(\nu,e)+\nabla u(0)\cdot(A\nu)=D^2u(0)(\nu,e),
\]
because $\nabla u(0)=\rho\nabla r(0)$ is normal and $A\nu=v$ is tangent. Hence $|D^2u(0)(\nu,e)|\le C$.

\emph{The normal direction.} By \eqref{eq:levi} with $u=\rho r$, at $0$ the compression of $H(u)$ to the
quaternionic orthogonal complement $\nu^\perp$ is $\rho H(r)|_{\nu^\perp}\ge b'c_rI$. Since $H(B_0)\ge0$,
\begin{equation}\label{eq:Dlow}
Q|_{\nu^\perp}\ge b'c_rI .
\end{equation}
Choose a real orthonormal basis $\eps_1=\nu,\eps_2,\dots,\eps_{4m}$ of $\R^{4m}$ and write
$M=D^2\psi(0)=\sum_{i,j}M_{ij}\eps_i\eps_j^{\rm t}$. By linearity and polarisation of
$H(\xi\xi^{\rm t})=\hat\xi\hat\xi^*$,
\[
Q=H(M)=M_{11}\nu\nu^*+R,\qquad R=\sum_{(i,j)\ne(1,1)}M_{ij}\,\tfrac12\bigl(\hat\eps_i\hat\eps_j^*+\hat\eps_j\hat\eps_i^*\bigr),
\]
and the entries $M_{ij}$ with $(i,j)\ne(1,1)$ are bounded by the two previous steps, so $|R|\le C$.
Let $f_2,\dots,f_m$ be a quaternionic orthonormal basis of $\nu^\perp$. In the basis $\nu,f_2,\dots,f_m$,
\[
Q=\begin{pmatrix}q_{11}&z^*\\z&D\end{pmatrix},\qquad q_{11}=M_{11}+\nu^*R\nu,\quad
z_i=f_i^*R\nu,\quad D=Q|_{\nu^\perp}=R|_{\nu^\perp},
\]
so $z$ and $D$ are bounded and $D\ge b'c_rI$ by \eqref{eq:Dlow}. If $m>1$, \eqref{eq:schur} gives
$g_t=\det Q=(\det D)(q_{11}-z^*D^{-1}z)$, so
\[
0<q_{11}=\frac{g_t}{\det D}+z^*D^{-1}z\le\frac{C_0}{(b'c_r)^{m-1}}+\frac{|z|^2}{b'c_r}\le C;
\]
if $m=1$, $q_{11}=Q=g_t\le C_0$. In both cases $|M_{11}|\le|q_{11}|+|R|\le C$, so $|D^2\psi(0)|\le C$.

\subsection{H\"older estimates and higher derivatives}
By \eqref{eq:hessreduce} and the boundary bound, $\psi_{ee}\le C$ for every unit $e$. Since
$\Delta\psi=\Re\tr Q>0$ by Lemma \ref{lem:hessian}, every eigenvalue of $D^2\psi$ is at least
$-(4m-1)C$. So $|D^2\psi|\le C$ on $\overline{\tilde\Omega}$, hence $|Q|\le C$; with $\det Q=g_t\ge c_0$ every
eigenvalue of $Q$ lies in an interval $[c,C]\subset(0,\infty)$. Thus $D^2\psi$ takes values in the compact
convex set
\[
\mathcal K=\{M:|M|\le C,\ cI\le H(M)\le CI\},
\]
on which $F$ is smooth and concave, and by \eqref{eq:ellipticity}
$DF(M)[\xi\xi^{\rm t}]=\hat\xi^*H(M)^{-1}\hat\xi\in[C^{-1}|\xi|^2,c^{-1}|\xi|^2]$ for $M\in\mathcal K$. Multiply $F$ by a
smooth cut-off supported in $\{H(M)>0\}$ and equal to one near $\mathcal K$, and extend by zero; the
equation $F(D^2\psi)=\log g_t$ is then a fully nonlinear, uniformly elliptic, concave equation on the
range of $D^2\psi$, with smooth dependence on $x$ through $g_t$ and boundary values $B_0$. These are
the hypotheses (i)$'$, (ii)$'$ of the interior and boundary estimates \cite[Theorems 17.14 and
17.26$'$]{GT01}, which are due to Evans \cite{Evans82} and Krylov \cite{Kry83,Kry84} and are local; in
the flat charts of a finite cover they give $\|\psi\|_{C^{2,\alpha}(\overline{\tilde\Omega})}\le C$ for some
$\alpha\in(0,1)$.

For higher derivatives we argue by induction: suppose $\|u\|_{C^{j+2,\alpha}}\le C_j$. In the interior,
for each constant direction $e$, $\psi_e$ solves the linear equation \eqref{eq:diffeq} whose
coefficients $Q^{-1}$ are bounded in $C^{j,\alpha}$ and uniformly elliptic; the interior Schauder
estimates bound $\psi_e$ in $C^{j+2,\alpha}$ on compact subsets of the chart. Near the boundary, flatten
it by a smooth change of coordinates $y=\Phi(x)$ with the boundary at $y_{4m}=0$. The equation becomes
$\mathcal G(y,D_yv,D_y^2v)=0$ for $v=u\circ\Phi^{-1}$, with $\mathcal G$ smooth and
$\partial\mathcal G/\partial v_{ij}$ uniformly elliptic on the range of the solution. For a tangential
coordinate direction $y_\iota$, $\iota<4m$, the derivative $v_\iota$ vanishes on $y_{4m}=0$, because
$v$ does, and solves the linear equation obtained by differentiating, with coefficients
$\partial\mathcal G/\partial v_{ij},\partial\mathcal G/\partial v_i$ bounded in $C^{j,\alpha}$ and right-hand side
$-\partial_{y_\iota}\mathcal G$ bounded in $C^{j,\alpha}$. The local boundary Schauder estimates bound
$v_\iota$ in $C^{j+2,\alpha}$ near the boundary. This controls in $C^{j+1,\alpha}$ every second derivative
$v_{il}$ except $v_{4m\,4m}$. Uniform ellipticity gives $\partial\mathcal G/\partial v_{4m\,4m}\ge c>0$, so
by the implicit function theorem $v_{4m\,4m}=\Psi(y,D_yv,(v_{il})_{(i,l)\ne(4m,4m)})$ with $\Psi$
smooth, and the right side is bounded in $C^{j+1,\alpha}$. Hence $\|u\|_{C^{j+3,\alpha}}\le C_{j+1}$. All
constants are independent of $t$, which completes the a priori estimates and the existence proof.

\subsection{Parameters}
Let the data depend smoothly on a parameter $\tau$ in an open set of $\R^K$, and let
$\Theta_\tau$ identify $\overline V_{\tau_0}$ with $\overline V_\tau$ as above. The gap \eqref{eq:gap} at
$t=1$ and the bounds $c_r$, $C_0$ are locally uniform in $\tau$ by continuity, so $b,b'$ can be fixed
locally uniformly. Pulling the equation back by $\mathrm{id}\times\Theta_\tau$ gives on the fixed domain
$\tilde\Omega_{\tau_0}$ an equation $\mathcal F_\tau(v)=0$, $v\in C^{k+2,\alpha}_0$, depending smoothly on
$(\tau,v)$. At $\tau_0$ its linearisation is $L_{\psi_{\tau_0}}$ conjugated by the pull-back, an
isomorphism by Lemma \ref{lem:lineardirichlet}. The implicit function theorem gives solutions
$v_\tau$ depending smoothly on $\tau$ with values in $C^{k+2,\alpha}_0$, for every $k$; by uniqueness they
are the pull-backs of the solutions $u_\tau$. A map from parameters to $C^{k+2,\alpha}$ that is smooth for
every $k$ is jointly smooth in $(\tau,x)$, which is the asserted dependence. The descent from
$\tilde\Omega$ to $\Omega$ is compatible with all of this.

Lemma \ref{lem:comparison2} also holds in three slots on subdomains of $\Omega$, in particular on the
balls of the fibres used in Section \ref{sec:fiber}. \qed
